\section{Two Sum IV: Input is a BST}
\label{sec:bst-two-sum}

\problemstatement{Given the root of a BST and a target integer $k$,
determine whether there exist two \textbf{distinct} nodes whose values
sum to exactly $k$.}

\begin{patternbox}
\emph{``Two sum''} on a BST is the classic two-pointer two-sum technique
(normally run on a sorted array) combined with
Section~\ref{sec:bst-iterator}'s BST Iterator: since an inorder traversal
already produces values in sorted order on demand, we can drive
\textbf{two} iterators -- one walking forward (ascending) from the
smallest value, one walking backward (descending) from the largest --
exactly as the classic two-pointer two-sum algorithm drives two pointers
inward from the two ends of a sorted array, without ever materializing
that sorted array.
\end{patternbox}

\subsection*{Understanding the problem carefully}

If the tree's values were laid out as a sorted array, the standard
two-sum technique would maintain a pointer at each end: compute the sum
of the two pointed-at values; if it equals $k$, done; if it is too
small, advance the left pointer rightward (to a larger value); if too
large, move the right pointer leftward (to a smaller value). A BST
supports exactly this access pattern \emph{without} an explicit array,
using two independent, resumable inorder-traversal iterators: one
producing ascending values on demand (Section~\ref{sec:bst-iterator}'s
iterator exactly as built), and a mirror-image iterator producing
descending values on demand (the same idea, with left/right swapped).

\begin{ideabox}[Two BST Iterators, Driven Like Two-Sum's Two Pointers]
Initialize a forward iterator $\textit{lo}$ (smallest-first) and a
backward iterator $\textit{hi}$ (largest-first), both via the stack-based
technique from Section~\ref{sec:bst-iterator}. Maintain their current
values $a = \textit{lo}.\textit{next}()$ and $b =
\textit{hi}.\textit{next}()$. While $a < b$: if $a+b=k$, return true; if
$a+b<k$, advance $\textit{lo}$ (get a larger $a$); if $a+b>k$, advance
$\textit{hi}$ (get a smaller $b$). If the pointers meet or cross ($a \ge
b$) without finding a match, return false.
\end{ideabox}

\subsection*{Why driving two independent stack-based iterators is correct and efficient}

This is exactly the two-pointer two-sum algorithm, with the ``array'' it
normally operates on replaced by two on-demand sorted streams generated
directly from the tree. Each iterator independently maintains its own
$O(h)$-space stack, advancing forward or backward through the BST's
sorted order exactly as Section~\ref{sec:bst-iterator} established, and
because both iterators only ever move \emph{toward} each other (never
revisiting a value), the total work across the entire search is bounded
by the number of \emph{distinct} values each iterator visits before they
cross -- giving the same linear-total-work guarantee as the classic
array-based two-pointer technique, without ever needing $O(n)$ space to
store a flattened sorted array.

\subsection*{Algorithm}

\begin{enumerate}
  \item Initialize a forward BST iterator $\textit{lo}$ and a backward
        BST iterator $\textit{hi}$.
  \item $a \gets \textit{lo}.\textit{next}()$, $b \gets
        \textit{hi}.\textit{next}()$.
  \item While $a < b$:
  \begin{enumerate}
    \item If $a+b=k$: return true.
    \item If $a+b<k$: $a \gets \textit{lo}.\textit{next}()$.
    \item Else: $b \gets \textit{hi}.\textit{next}()$.
  \end{enumerate}
  \item Return false.
\end{enumerate}

\subsection*{Dry run}

\dryrun{Take the BST with root $5$, left child $3$, right child $8$,
target $k=10$.}

Forward iterator yields $3,5,8$ in order; backward iterator yields
$8,5,3$ in order.

\begin{itemize}
  \item $a=\textit{lo.next}()=3$, $b=\textit{hi.next}()=8$. $a<b$.
        $a+b=11>10$. Advance \textit{hi}: $b=5$.
  \item $a=3<b=5$. $a+b=8<10$. Advance \textit{lo}: $a=5$.
  \item $a=5$, $b=5$: $a<b$ is now false ($5 \not< 5$). Loop ends.
\end{itemize}

Return false. Checking by hand: the only pair sums available are
$3{+}5{=}8$, $3{+}8{=}11$, $5{+}8{=}13$ -- none equal $10$. Correct.

\subsection*{C++ implementation}

\begin{lstlisting}
#include <bits/stdc++.h>
using namespace std;

struct TreeNode {
    int val;
    TreeNode* left;
    TreeNode* right;
    TreeNode(int v) : val(v), left(nullptr), right(nullptr) {}
};

class BSTIterator {
    stack<TreeNode*> st;
    bool reverseOrder; // true = descending (successor uses right first)

    void pushChain(TreeNode* node) {
        while (node) {
            st.push(node);
            node = reverseOrder ? node->right : node->left;
        }
    }

public:
    BSTIterator(TreeNode* root, bool reverse) : reverseOrder(reverse) {
        pushChain(root);
    }

    int next() {
        TreeNode* node = st.top();
        st.pop();
        pushChain(reverseOrder ? node->left : node->right);
        return node->val;
    }
};

bool findTarget(TreeNode* root, int k) {
    BSTIterator lo(root, false); // ascending
    BSTIterator hi(root, true);  // descending

    int a = lo.next(), b = hi.next();
    while (a < b) {
        if (a + b == k) return true;
        else if (a + b < k) a = lo.next();
        else b = hi.next();
    }
    return false;
}

int main() {
    TreeNode* root = new TreeNode(5);
    root->left = new TreeNode(3);
    root->right = new TreeNode(8);

    cout << (findTarget(root, 10) ? "true" : "false") << endl;
    cout << (findTarget(root, 11) ? "true" : "false") << endl;
    return 0;
}
\end{lstlisting}

\begin{complexitybox}
\textbf{Time Complexity.} $O(n)$ in the worst case -- each value is
produced by one of the two iterators at most once before they cross,
and each iterator's amortized cost per call is $O(1)$, as established in
Section~\ref{sec:bst-iterator}.
\textbf{Space Complexity.} $O(h)$ for each iterator's stack, so $O(h)$
total.
\end{complexitybox}

\begin{takeawaybox}
A BST is, for the purposes of any algorithm that only needs sorted-order
access, a drop-in replacement for a sorted array -- as long as you have
an $O(h)$-space, resumable iterator (forward and/or backward) instead of
an $O(n)$-space materialized array. Any two-pointer array technique
(two-sum, partitioning, merging) can be ported to a BST this way, exactly
as done here.
\end{takeawaybox}
